% GENERATED FILE. Edit canonical lessons, then run npm run generate:books.
% canonical-source-hash: fnv1a64:b82be7793e8fbb8e
% canonical-lessons: PT-C307-ferido, PT-C307-ferir, PT-C307-ferimento, PT-C307-fratura, PT-C307-gesso

\chapter{Feridos}
\label{ch:pt-feridos}
\hlchaptermodality{307}

\begin{chapteropening}
\textbf{By the end of this chapter:} \emph{I can talk about the injured, their injuries and fractures, and a plaster cast.}

Complete the last lesson of chapter 307: I can talk about the injured, their injuries and fractures, and a plaster cast.
\end{chapteropening}

\section[o ferido]{o ferido — the injured person, the wounded}

\label{lesson:PT-C307-ferido}



\begin{quote}
Before the new one: say the Portuguese for to survive, then the Portuguese for the survivor.
\end{quote}

\subsection*{You'll want to know: o ferido}
\textbf{o ferido} — \textquotedblleft{}the injured person, the wounded\textquotedblright{}.

\begin{quote}
O ferido foi levado para o hospital.
\end{quote}

\begin{quote}
\emph{The injured man was taken to hospital.}
\end{quote}

\emph{Feridos ligeiros} — minor injuries; \emph{feridos graves} — serious ones.

\begin{cousinweb}
\emph{ferido} comes from Latin \emph{ferītus} \textquotedblleft{}struck\textquotedblright{} (EtymDB), from \emph{ferīre} \textquotedblleft{}to strike\textquotedblright{}.
\end{cousinweb}

\begin{grammarlens}[title={What you've built}]
One more word for telling, explaining and describing at B1.
\end{grammarlens}

\subsection*{Your turn}
\begin{itemize}
\raggedright
  \item \emph{Say it:} \emph{o ferido}
  \item \emph{Say it:} \emph{o ferido}, once more
  \item \emph{Say it:} a sentence of your own with \emph{o ferido}
  \item \emph{Recall:} say the Portuguese for to survive, then the Portuguese for the survivor, then say \emph{o ferido} again
\end{itemize}

\begin{tcolorbox}[breakable,colback=teal!4,colframe=teal!35!black,title={Before you move on}]
What does o ferido mean? (\textquotedblleft{}The injured person, the wounded\textquotedblright{}.) Say it once more.
\end{tcolorbox}
\section[ferir]{ferir — to injure, to wound; to hurt}

\label{lesson:PT-C307-ferir}



\begin{quote}
Before the new one: say the Portuguese for the survivor, then the Portuguese for the injured person.
\end{quote}

\subsection*{You'll want to know: ferir}
\textbf{ferir} — \textquotedblleft{}to injure, to wound; to hurt\textquotedblright{}.

\begin{quote}
A queda feriu-lhe a perna.
\end{quote}

\begin{quote}
\emph{The fall injured his leg.}
\end{quote}

Present: \emph{firo, feres, fere}. Also feelings: \emph{ferir alguém}.

\begin{cousinweb}
\emph{ferir} is inherited from Latin \emph{ferīre} \textquotedblleft{}to strike\textquotedblright{} (EtymDB).
\end{cousinweb}

\begin{grammarlens}[title={What you've built}]
One more word for telling, explaining and describing at B1.
\end{grammarlens}

\subsection*{Your turn}
\begin{itemize}
\raggedright
  \item \emph{Say it:} \emph{ferir}
  \item \emph{Say it:} \emph{ferir}, once more
  \item \emph{Say it:} a sentence of your own with \emph{ferir}
  \item \emph{Recall:} say the Portuguese for the survivor, then the Portuguese for the injured person, then say \emph{ferir} again
\end{itemize}

\begin{tcolorbox}[breakable,colback=teal!4,colframe=teal!35!black,title={Before you move on}]
What does ferir mean? (\textquotedblleft{}To injure, to wound; to hurt\textquotedblright{}.) Say it once more.
\end{tcolorbox}
\section[o ferimento]{o ferimento — the injury, the wound}

\label{lesson:PT-C307-ferimento}



\begin{quote}
Before the new one: say the Portuguese for the injured person, then the Portuguese for to injure.
\end{quote}

\subsection*{You'll want to know: o ferimento}
\textbf{o ferimento} — \textquotedblleft{}the injury, the wound\textquotedblright{}.

\begin{quote}
O ferimento na cabeça não é grave.
\end{quote}

\begin{quote}
\emph{The head injury is not serious.}
\end{quote}

Built on \emph{ferir}.

\begin{cousinweb}
\emph{ferimento} is built on \emph{ferir} \textquotedblleft{}to injure\textquotedblright{}, from Latin \emph{ferīre} \textquotedblleft{}to strike\textquotedblright{} (EtymDB), with the noun ending \emph{-mento}.
\end{cousinweb}

\begin{grammarlens}[title={What you've built}]
One more word for telling, explaining and describing at B1.
\end{grammarlens}

\subsection*{Your turn}
\begin{itemize}
\raggedright
  \item \emph{Say it:} \emph{o ferimento}
  \item \emph{Say it:} \emph{o ferimento}, once more
  \item \emph{Say it:} a sentence of your own with \emph{o ferimento}
  \item \emph{Recall:} say the Portuguese for the injured person, then the Portuguese for to injure, then say \emph{o ferimento} again
\end{itemize}

\begin{tcolorbox}[breakable,colback=teal!4,colframe=teal!35!black,title={Before you move on}]
What does o ferimento mean? (\textquotedblleft{}The injury, the wound\textquotedblright{}.) Say it once more.
\end{tcolorbox}
\section[a fratura]{a fratura — the fracture, the break}

\label{lesson:PT-C307-fratura}



\begin{quote}
Before the new one: say the Portuguese for to injure, then the Portuguese for the injury.
\end{quote}

\subsection*{You'll want to know: a fratura}
\textbf{a fratura} — \textquotedblleft{}the fracture, the break\textquotedblright{}.

\begin{quote}
Ela tem uma fratura no braço.
\end{quote}

\begin{quote}
\emph{She has a fractured arm.}
\end{quote}

Before the 1990 spelling reform, \emph{fractura}.

\begin{cousinweb}
\emph{fratura} is a learned borrowing from Latin \emph{frāctūra} \textquotedblleft{}a breaking\textquotedblright{}, from \emph{frangere} \textquotedblleft{}to break\textquotedblright{} (Lewis \& Short).
\end{cousinweb}

\begin{grammarlens}[title={What you've built}]
One more word for telling, explaining and describing at B1.
\end{grammarlens}

\subsection*{Your turn}
\begin{itemize}
\raggedright
  \item \emph{Say it:} \emph{a fratura}
  \item \emph{Say it:} \emph{a fratura}, once more
  \item \emph{Say it:} a sentence of your own with \emph{a fratura}
  \item \emph{Recall:} say the Portuguese for to injure, then the Portuguese for the injury, then say \emph{a fratura} again
\end{itemize}

\begin{tcolorbox}[breakable,colback=teal!4,colframe=teal!35!black,title={Before you move on}]
What does a fratura mean? (\textquotedblleft{}The fracture, the break\textquotedblright{}.) Say it once more.
\end{tcolorbox}
\section[o gesso]{o gesso — the plaster (cast)}

\label{lesson:PT-C307-gesso}



\begin{quote}
Before the new one: say the Portuguese for the injury, then the Portuguese for the fracture.
\end{quote}

\subsection*{You'll want to know: o gesso}
\textbf{o gesso} — \textquotedblleft{}the plaster (cast)\textquotedblright{}.

\begin{quote}
O meu irmão anda de gesso há um mês.
\end{quote}

\begin{quote}
\emph{My brother has been in a plaster cast for a month.}
\end{quote}

Also plaster for walls.

\begin{cousinweb}
\emph{gesso} comes from Latin \emph{gypsum} (EtymDB), from Greek \emph{gýpsos} \textquotedblleft{}chalk, gypsum\textquotedblright{}.
\end{cousinweb}

\begin{grammarlens}[title={What you've built}]
One more word for telling, explaining and describing at B1.
\end{grammarlens}

\subsection*{Your turn}
\begin{itemize}
\raggedright
  \item \emph{Say it:} \emph{o gesso}
  \item \emph{Say it:} \emph{o gesso}, once more
  \item \emph{Say it:} a sentence of your own with \emph{o gesso}
  \item \emph{Recall:} say the Portuguese for the injury, then the Portuguese for the fracture, then say \emph{o gesso} again
\end{itemize}

\begin{tcolorbox}[breakable,colback=teal!4,colframe=teal!35!black,title={Before you move on}]
What does o gesso mean? (\textquotedblleft{}The plaster (cast)\textquotedblright{}.) Say it once more.
\end{tcolorbox}
